% Existing static assets only: one PDF per stem, including every gallery page.
% Nonfloating blocks keep each caption with its image and avoid float-only pages.
\clearpage
\section{Figures and numerical illustrations}
\label{cp:figures}
\begingroup
\raggedbottom
\captionsetup{font=small,labelfont=bf,skip=4pt,hypcap=false}
\newcommand{\figureplate}[4][1]{%
  \par\noindent\begin{minipage}{\linewidth}
  \centering
  \includegraphics[page=#1,width=\linewidth,height=#2\textheight,
                   keepaspectratio]{figures/#3.pdf}
  \captionof{figure}{#4}\label{fig:archive:#3:#1}
  \end{minipage}\par\medskip}

This collection preserves the static figures already accompanying the paper,
including alternate layouts and all seven pages of the illustrated gallery.
Format copies and interactive explainers are omitted. These are earlier
illustrations under their recorded protocols, not new experiments or
certificates of the headline theorem. Unless stated otherwise, RMS means
disagreement with dense predictions on sampled queries, not error against
labels. Recorded times do not certify a continuous-time supremum, and
separately fitted endpoints need not share a physical time. Historical
artwork sometimes writes $P$ for the history order $q$. In radial plots,
$f(\theta)$ denotes the prediction at input angle $\theta$.

\figureplate{.23}{mechanism}{Response-memory reconstruction and the
product-of-tails identity. Both histories use the same Legendre modes;
the initialized mixer remains fixed and retained. The printed finite-horizon
rates are schematic, not the present complete-trajectory estimates.}
\figureplate{.32}{memory_mechanism}{Single-sample, three-mode illustration:
paired forward and residual-weighted backward moments reconstruct the learned
connection. The reconstructed network supplies the next responses; the
polynomial basis is fixed while its neuron-valued coefficients evolve.}

\figureplate{.31}{paired_projection}{Projecting both histories onto the same
orthogonal history space cancels the two mixed terms, leaving the pairing of
the omitted histories. Trajectory control additionally requires regularity
and feedback stability.}
\figureplate{.25}{clocks}{Illustrative four-mode fits of a normalized backward
history in physical time and two alternative clocks. The displayed fit errors
describe this example, not a general ordering of clock accuracy.}
\figureplate{.19}{moments}{Three neuron histories reconstructed with one,
two, and three Legendre coefficients, with coefficient distributions below.
Errors are normalized by each history's range. This older nonzero-readout
example illustrates memory coordinates, not the canonical initialization
theorem.}

\figureplate{.39}{learning_clocks}{A constructed smooth response path in
physical time and a clock that also tracks response motion. Matching colors
identify the same events; this is a coordinate-change schematic, not a
measured training trajectory or a constant-speed assertion.}
\figureplate{.39}{learning_in_motion}{Earlier eight-checkpoint circle
comparison: two tanh hidden layers, $n=2048$, and $q=1,3,7$.
Snapshots share physical times; discrepancies use 2048 query angles.
Dotted numerical sensitivities and connecting lines do not certify the
unobserved times.}

\figureplate{.20}{trajectory}{Paired-label circle predictions at common times
$0,5,20,80$: two tanh hidden layers, $n=2048$, $q=1,3,7$.
Radius is $3+f(\theta)$ and black points are training labels.
This is the later 39-checkpoint, two-tolerance Heun replay.}
\figureplate{.28}{trajectory_rms}{Companion dense-versus-memory RMS at the
38 positive common times of that replay. The coarse/fine sensitivity scale
is an empirical numerical diagnostic, not a continuous-time bound.}
\figureplate{.24}{orders}{Deep-circle endpoint discrepancies, coarse/fine
sensitivity, and hidden-feature motion for $q=1,2,3$.
These three-hidden-layer tanh runs have $n=4096$ and individual training-MSE
$10^{-3}$ endpoints; the sensitivity ranges are not confidence intervals.}

\figureplate{.59}{circle_deep}{Five circle tasks with three tanh hidden
layers, $n=4096$, and $q=1,2,3$. Each curve uses its own training-MSE
$10^{-3}$ endpoint: the comparison is between fitted functions, not common
physical times.}
\figureplate{.23}{circles_deep}{Compact radial view of the deep-circle
family. Radius is $3+f(\theta)$; the printed values are RMS differences from
dense predictions on 8192 circle queries. This is an alternate layout,
not an additional experiment.}

\figureplate{.43}{circle_shallow}{Earlier Cartesian circle comparison with
two tanh hidden layers, $n=2048$, and $q=1,3,7$, at individual training-MSE
$10^{-3}$ endpoints. This export uses archived dense references, unlike
the selected references in the radial views.}
\figureplate{.40}{circle_shallow_radial}{Shallow-circle radial comparison at
the same loss threshold. Radius is $3+f(\theta)$; RMS uses 8192 queries
against four fresh dense references and one refined archived reference.
Black points mark training labels.}

\figureplate{.59}{circle_deep_radial}{Larger radial panels for the deep-circle
family: three tanh hidden layers, $n=4096$, $q=1,2,3$, and separately fitted
training-MSE $10^{-3}$ endpoints. The printed discrepancies use 8192 circle
queries; fixed dense mixers remain additional storage.}
\figureplate{.23}{circles_shallow}{Compact radial layout of the shallow-circle
family, using the same selected references as its larger radial view.
The printed RMS values measure fitted-function disagreement, not label risk.}

\figureplate{.39}{order_decay}{Observed circle discrepancy versus history
order in five tasks. The inverse-order reference lines are visual comparisons,
not fitted or certified convergence rates.}

\figureplate{.25}{same_rank}{Selected circle endpoints compare response
memory with two directly trained factor seeds at the same correction-rank
bound. Equal correction rank does not imply equal total storage or a
comparison with every low-rank optimizer.}
\figureplate{.39}{factors}{Paired-label fitted functions and five-task rank
sweeps for two tanh hidden layers of width 2048. RMS uses 8192 queries at
individual training-MSE $10^{-3}$ endpoints; fixed mixers are additional,
and omitted or precision-limited runs are not exact successes.}

\figureplate{.39}{same_rank_dynamics}{Alternate presentation of the factor
comparison, including the rank-24 paired-label example and five-task sweep.
Despite the filename, these are separately fitted endpoints, not common-time
dynamics. Hollow points have limited numerical precision; one capped run is
omitted.}
\figureplate{.39}{same_rank_radial_preview}{Enlarged paired-label comparison:
dense training, order-3 memory, and two factor seeds with correction-rank
bound 24. Radius is $3+f(\theta)$; RMS uses 8192 angles at individual
training-MSE $10^{-3}$ endpoints.}

\figureplate{.39}{frozen_ntk_radial_preview}{Dense training, its full frozen
initial empirical NTK, and order-3 memory from the same width-2048
initialization. Each reaches training MSE $10^{-3}$ at its own time.
Radius is $30+f(\theta)$ here, preserving the large signed NTK excursions.}
\figureplate{.36}{learning_controls_equal_mixed_odd}{Equally spaced
mixed-label control: dense training, frozen initial NTK, two factor seeds,
and order-3 memory. RMS uses 8192 queries at individual training-MSE
$10^{-3}$ endpoints. The correction-rank bound is 12 after the antipodal
quotient.}

\figureplate{.36}{learning_controls_quadrant_alternating}{Quadrant
alternating labels, with the same width-2048 two-hidden-layer tanh control
protocol. The factor correction-rank bound is 24; the NTK panel explicitly
uses a 100-times-wider output scale. RMS is disagreement with dense,
not ground-truth error.}
\figureplate{.36}{learning_controls_quadrant_center_edges}{Quadrant
center-versus-edges labels under the same endpoint protocol. The NTK panel
alone uses a 200-times-wider output scale, with signed ticks in original
units. All reported RMS discrepancies use 8192 circle queries.}

\figureplate{.36}{learning_controls_quadrant_pairs}{Quadrant paired labels:
the same dense, frozen-NTK, rank-24 factor, and order-3 memory comparison.
Each model fits separately to training MSE $10^{-3}$; the NTK output scale
is explicitly widened tenfold.}

\figureplate{.82}{learning_controls_gallery}{Overview of all five preceding
control tasks. This combines the same observations, with marked NTK scale
changes and correction-rank bounds 24 or 12; it is not another replicate.}
\figureplate{.51}{learning_controls_gallery_a}{Page-friendly assembly of
the three quadrant control tasks. The protocols, separate fitted endpoints,
and marked NTK output-scale changes are unchanged.}
\figureplate{.34}{learning_controls_gallery_b}{Companion assembly of the
outlier and equally spaced mixed-label tasks. These are the remaining two
rows of the full control gallery, not additional experiments.}

\figureplate{.27}{mnist}{Earlier MNIST 3-versus-8 comparison: 100 training
and 1984 held-out images, two tanh hidden layers, $n=4096$.
The order-3 agreement plot and order-1--3 differences use individual
training-MSE $10^{-3}$ endpoints. RMS is disagreement with dense predictions.}
\figureplate{.46}{mnist_rms}{MNIST validation discrepancy across history
orders and the displayed training-loss endpoints. Bars show observed
refinement margins, not confidence intervals; open markers identify failures
of the numerical gate.}

\figureplate{.39}{mnist_residuals}{Alternate MNIST endpoint presentation.
Signed memory-minus-dense differences share one vertical scale; order 2 is
slightly closer than order 3 in this run. These are function-agreement
measurements, not classification accuracies.}
\figureplate{.25}{mnist_scatter}{Held-out MNIST predictions for $q=1,2,3$
against dense predictions at the primary training-MSE $10^{-3}$ endpoints.
This scatter layout uses the same endpoint data, not another replicate.}

\figureplate{.41}{sphere_orders}{Sphere illustration with four ReLU hidden
layers, $n=2048$, $m=64$, and 8192 equal-area queries. Both hemispheres and
a shared error scale are shown. Each model stops at training RMS $0.01$;
the single-seed Euler run uses step $1/128$. ReLU is outside the analytic
activation theorem.}
\figureplate{.41}{sphere_comparison}{Dense output, order-3 memory output,
and their signed difference for the same sphere run. Black points are
training locations. This fitted-endpoint illustration has no
continuous-flow refinement certificate.}

\figureplate{.39}{sphere3_time}{Saved sphere-3 query RMS from a width-4096
dense reference over recorded times, with eight training and 30 test inputs.
The export's ``Logarithmic'' family is called Taylor by its adjacent plotting
script, not the seeded decoder below. Dense-control curves aggregate three
small-network seeds and two dense pairs by geometric means.}
\figureplate{.39}{sphere3_tradeoff_endpoint}{Sphere-3 fitted-endpoint RMS
versus moving parameter count. Smaller-dense summaries are medians and
observed ranges over three seeds. ``Logarithmic'' retains the same historical
Taylor-family label; fixed storage and offline setup are additional.}

\figureplate{.39}{sphere3_tradeoff_max_time}{Companion sphere-3 tradeoff for
maximum RMS over recorded times, rather than endpoint RMS. The shaded
small-network range is not a confidence interval; the horizontal axis still
excludes fixed storage and source-construction costs.}
\figureplate{.26}{compression_validation}{Earlier Harmonic circle
confirmation at widths 1024, 2048, and 4096, with three held-out seeds per
width. Curves show maximum-recorded-time query RMS and retained tensor
scalars; shading is the observed seed range. Source setup and query workspace
are additional.}

\figureplate{.26}{compression_logarithmic}{Earlier dimension-five
\emph{seeded decoder}, labelled ``Logarithmic'' in this artwork: finite-program
query fidelity and persistent numerical payload across three widths.
This is not Taylor. The dotted storage curve adds query workspace; flat
payload in these tests is not an asymptotic scaling estimate.}
\figureplate{.23}{compression_checks}{Further Harmonic dimension and
PCA-based digits checks, followed by seeded-decoder dimension checks.
The decoder's circle failure remains visible. Seed counts vary; ranges and
the empirical factor-three threshold are not confidence statements.
The digits task differs from the MNIST100 figures above.}
\figureplate{.26}{learning_controls_two_outliers_alternating}{Remaining
circle control: alternating labels with two outliers, under the earlier
width-2048 tanh endpoint protocol. The NTK output scale is widened tenfold;
RMS measures fitted-function disagreement on 8192 queries.}

\figureplate[1]{.39}{experimental_gallery}{Original illustrated gallery,
page 1: the three-mode memory mechanism. The next six plates preserve the
remaining pages, including their explanatory headings; they reuse earlier
illustrations rather than add experiments.}
\figureplate[2]{.39}{experimental_gallery}{Gallery page 2: the earlier
eight-checkpoint common-time circle comparison, not the later 39-checkpoint
replay. Its numerical sensitivities are diagnostics between saved runs.}
\figureplate[3]{.39}{experimental_gallery}{Gallery page 3: same-correction-rank
factor controls at separately fitted endpoints. Rank equality is not an
equal-total-storage claim.}
\figureplate[4]{.39}{experimental_gallery}{Gallery page 4: the constructed
smooth-path clock illustration. The change of coordinate preserves events,
not necessarily their speed.}
\figureplate[5]{.39}{experimental_gallery}{Gallery page 5: orthogonal pairing
of two histories leaves a product of omitted tails. Feedback stability is a
separate requirement for controlling the evolving network.}
\figureplate[6]{.39}{experimental_gallery}{Gallery page 6: the earlier
MNIST100 fitted-prediction comparison and signed differences. RMS measures
disagreement with dense predictions, not error against labels.}
\figureplate[7]{.39}{experimental_gallery}{Gallery page 7: deep-circle order
sensitivity and hidden-feature motion. Refinement scales are empirical
diagnostics, and the observed order trend need not be monotone.}
\figureplate{.39}{order_sensitivity}{Standalone view of the deep-circle order,
numerical-sensitivity, and feature-motion comparison. Each model uses its own
training-MSE $10^{-3}$ endpoint; order 2 is closer than order 3 on the two
alternating-label tasks, so the observed order trend is not uniformly monotone.}
\endgroup
